\subsection{工作流引导、科学判断与展望}

主对照表明，Cloudra通过使化学对象、计算结果与科学要求在探索中持续衔接，能够提高同一backbone的研究交付质量。Hard任务中一致的改善及Normal任务的收益，将这一贡献定位到比较性立体化学研究和多步研究的推进。完整50题分解显示，证据闭合与目标准确度均有收益；Hard中两项加权贡献分别为+0.1353和+0.1274。HH1和HI1进一步说明，即使已得到接近参考的能垒，数值结果与支撑其科学归属的路径证明仍是两个不同的研究成果。

归档案例提示，自主研究过程需要保持三种能力。首先，科学对象的身份需要超越当前语言描述而持续存在。一个反应物种可以包含多个构象，所要求步骤的初猜也可能优化至属于另一条路线的鞍点。HA8需要将能量分母持续关联至稳定物种，同时保留IRC实际连接的局部端点；NU6需要按优化后的连接关系判断新驻点，而非从初猜的预期名称继承身份。这些关系跨越了不同时间得到的计算结果。持续保存的研究图使模型能够在每次计算改变已知事实后，重新检查并更新这些关系。 这一解释与语言模型研究中外部规划、检验及反馈的作用相呼应。

其次，反馈需要指出研究应在哪一层改变。缺少端点优化与鞍点反复落到另一条路线需要不同动作，两者也不等同于方法无法描述目标电子态。语言模型依据当前上下文生成合理的后续动作；熟悉的化学表述可以保持连贯，而数值或路径证据已经发生变化。外部测量与结构化证据能够暴露这种失配，但智能体仍需判断应当修复计算、修改路径假设、重新选择参照，还是拒绝给出目标结论。Record对照说明，使工具可用并不足以保证这些关系持续进入最终论证；Brief-errors对照则说明，反馈中什么信息有用，取决于当前需要作出的决定。 文本解释对既有判断的合理化也为这种语义与证据脱节提供了相关研究背景。

第三，完成局部证明与判断科学主张的范围需要共同推进。Brief-errors较高的Impossible均分来自IA1和IA9的特定目标结论，IG8则呈相反结果。其IC3图提供了更具体的行为证据：运行保留已表征的闭壳层分支，同时拒绝将其当作所求最低开壳层单重态，或把其相对能量报告为目标能隙。这是一种有价值的科学判断：保留探索已经建立的结果，同时不让方便得到的结果替换原本不同的研究问题。交叉结果支持一个依赖任务的解释假说：详细诊断有助于补齐有效的局部证明，但当方法或目标本身需要重新审视时，持续修补当前证明可能无助于解决问题。本次最终结果对照并未确立诊断长度的因果作用，也未测量其背后的模型内部认知机制。

分数与主张记录还显示了评价器需要进一步区分的科学身份。Record/IA5正确拒绝目标异旋路径，但归档裁定记录了匹配器选择另一条相同端点正路径的情况；Brief-errors/IC3尽管明确区分目标开壳层态及独立闭壳层分支，准确度仍为0。本文保留定稿分数，并与保存的科学主张一同解释。这些案例说明，仅有分子连接关系不足以标识反应通道或电子态目标；评价过程也需要保留研究工作流旨在维持的状态与路径区别。

P2每题每配置保留一次运行，任务簇bootstrap区间反映题目构成变化，而非同题独立重复的运行方差。完整整理表覆盖全部200次运行及其三项评分分量。部分Hard任务的准确度目标包含定性通道身份，因此所报告的收益包括研究完成质量与数值能垒。后续在相同任务上重复评估，可以进一步刻画运行间变异，以及不同研究路线在何种条件下被选择。

Impossible组按允许方法契约定义，不是对相应化学普遍不可能发生的断言。其分量还采用负目标规则，显式负结论的自洽性1与正路径的TS—IRC—端点证明分别解释。14例自洽性1、准确度0因而需要结合目标主张与来源理解。NU6能量指纹吻合则定位到可达路径中的相关问题：当鞍点身份与所要求步骤没有持续关联时，同一个能差可能被赋予不同科学含义。这些案例支持在研究假设改变时，持续检查每次计算实际回答了什么。

耗时与投入同样按实际观测单位报告。12个先导批次测量并发度2的四题真实耗时，包含原始失败与停止结果；P2投入档案覆盖选定的深挖研究，四配置共有15题。累计job-hours代表记录的量子化学进程时长，Standard/M1另有已注明的原生日志重构；它与核时、费用及P2端到端研究耗时分别记录。并发计算共享资源，CPU排队可能延迟个别结果。因此，这些测量支持已观察工作流与投入之间的比较；进一步区分排队、计算及智能体活跃时间，有助于明确实际耗时花在何处。

交付案例还给出了与科学改进相对应的系统方向。等待计算、收到完成事件、恢复研究及形成最终产物，需要共享连续的任务生命周期。Standard的缺图源于计算尚未完成时的提前结束；Guided/HI1也暴露了通知没有唤醒等待会话的问题。使未完成目标持续跨越这些状态变化，才能把新完成计算变成可用的研究信息。科学层面还需要让反复返回同一驻点区域的现象保持可见，帮助智能体判断再做一次局部尝试是否足以解决目标路径。更好的候选选择及研究假设修订，可以进一步补充证据保存能力的改善。

Cloudra缓解了语言模型直接进行几何推理与操作的困难，其构象采样能力相较早期版本也已有明显提升，但高柔性体系仍然具有挑战。随着构象多样性增加，寻找相关的反应物排布、中间体和反应路径可能需要更多采样与量子化学计算。提高构象探索效率，并在相互竞争的候选结构之间更合理地分配计算资源，是后续发展的重要方向。

我们的长期目标是将Cloudra从小分子机理计算扩展至更广泛的计算化学领域。多种计算引擎支持、共享执行基础设施及持续保存的科学记录，已为这一扩展提供基础；面向新的问题类型，还需要相应的科学表示、计算方法与验证准则。我们希望Cloudra最终成为计算化学家的强大研究工具，以及实验团队的计算专家。Cloudra将以GNU Affero General Public License（AGPL）开源，支持系统的公开检查、复用与社区共同发展。
